\section{Where the axioms come from}\label{sec:discussion}

This closing section is interpretation (\metag). It contains no theorem; each paragraph draws on a
result above, but the picture they compose is a reading, not a further result.

\subsection{Two limitative theorems, one pattern}

The foundations of mathematics carry two famous limits. G\"odel's theorems \citep{Godel1931} say
that a consistent, effectively axiomatized theory containing enough arithmetic does not prove
every arithmetical truth and, under the usual conditions on provability, does not prove its own
consistency. Cohen's theorems \citep{Cohen1963} say that natural questions, the continuum
hypothesis first among them, are not settled by the standard axioms. The two are usually filed
separately, as a limit on proof and a limit on decision, and usually as bad news.

The construction of this paper suggests filing them together. Its organizing patterns are that a
fixed rule, iterated, \emph{saturates}---it stays inside a closure it cannot leave---and that
content outside that closure comes only from \emph{extension}, the adjunction of something the
current package cannot express. Both patterns appear at the scale of
set theory. Saturation is Theorem~\ref{thm:nonattain}: the finite packaging rules never leave
$\HF$, which is why Infinity has to be purchased. Extension is the forcing analysis: a generic real
lies outside the ground model and strictly enlarges its algebra of coded predicates
(Theorems~\ref{thm:novelty} and \ref{thm:strict}).

On this reading G\"odel is the saturation face: a fixed, effectively given theory, run as long as
one likes, does not reach certain truths about itself. Cohen is the extension face: a ground model
is enlarged by adjoining a real it cannot code, and in suitable extensions old questions, such as
the size of the continuum, receive different answers. Two cautions keep the reading honest. The
G\"odel half needs all of its hypotheses: complete consistent theories exist, such as true
arithmetic, but they are not effectively axiomatized. And an independent \emph{sentence} is not the
same thing as a missing \emph{predicate}: Theorem~\ref{thm:strict} is about predicates on $\omega$,
and nothing here identifies the sentences a theory leaves undecided with the reals a model fails to
contain. The resemblance is one of shape---a closed apparatus run to saturation, and the lawful
extension beyond it---and it is offered as such.

\subsection{A foundation is an audited ledger}

If this is right, a foundation is neither a revealed truth nor an arbitrary convention, but a
\emph{ledger}. The distinction it draws, between questions settled inside a chosen framework and
the external question of which framework to adopt, is Carnap's \citep{Carnap1950}.
Table~\ref{tab:factorization} is that ledger written out for $\mathsf{ZFC}$. The free entries are
supplied by a modest packaging instrument and hold in $\HF$, with selected finite cores
machine-checked; at infinite stages they are carried by a structural premise that applies the same
clauses everywhere. The purchased entries are four named commitments, each shown to be
non-redundant, and the first boundary is sharp: $\HF$ satisfies everything except Infinity
(Proposition~\ref{prop:nonredundancy}(i)).

This is how $\mathsf{ZFC}$ can be a flat list and still be faithful to the staged construction:
over $\HF$ the staging is recoverable from membership alone (Theorem~\ref{thm:translation}). The
list is also not arbitrary, since its dependencies are fixed by what packaging supplies and what
must be bought. One can choose differently---a different lens (\textbf{CM-1}), no Foundation
(\textbf{CM-2}), no power-set totalization (\textbf{CM-7}), no selection principle
(\textbf{CM-4})---and the countermodel atlas is a catalogue of such alternative ledgers, in the
limit the plurality of set-theoretic universes discussed by \citet{Hamkins2012}. On this view
$\mathsf{ZFC}$ is one audited ledger among several, distinguished not by being true of an
antecedently given universe of objects \citep{Benacerraf1965} but by which commitments it buys.

\subsection{The open questions are expected}

On this reading the famous open questions lose some of their air of scandal. A foundation that
grows by lawful extension should have undecided questions: places where different extensions give
different answers. The continuum hypothesis is the standard example (\S\ref{sec:ch}). Over a
suitable ground model, one extension keeps it and another refutes it, and the ground's coded
predicates do not decide between them. Independence, read this way, is the expected signature of a
foundation that is still being extended, not a defect of one that was supposed to be finished.

\subsection{The answer}

We began with a list and a question: where does the list come from? Not from a pre-existing
universe of sets that the axioms approximately describe; we built no such universe and needed
none. Not from arbitrary convention either: the free entries are supplied by the chosen membership
instrument, the purchases are priced, and the order of dependence is not ours to choose. The list
comes from the place the framework locates stable objects in general: \emph{packaging under limited
access}---collecting multiplicities into units seen only through a lens---\emph{audited so that the
bookkeeping stays honest}, and \emph{extended when it saturates}. Set theory flattens that history
into a single membership relation; rank-erasure is the flattening; and forcing shows that the
history can always be continued.\footnote{The question ``where mathematics comes from'' is also
the title of \citet{LakoffNunez2000}, whose thesis concerns the cognitive origin of mathematical
ideas. That question is orthogonal to the structural account of where the \emph{axioms} of a
foundation come from offered here.}
